\documentclass[12pt]{article}
\usepackage{tcolorbox}
\usepackage{amsmath}

\include{preamble}
%All credit goes towards Dr. Adam Kapelner for providing the preamble and homework template.

\newtoggle{professormode}



\title{MATH 241 Spring \the\year~ Homework \#4}

\author{Steven Ou} %STUDENTS: write your name here

\iftoggle{professormode}{
\date{Due via Brightspace 11:59PM Sunday, March 8, \the\year \\ \vspace{0.5cm} \small (this document last updated \today ~at \currenttime)}
}

\renewcommand{\abstractname}{Instructions and Philosophy}

\begin{document}
\maketitle

\iftoggle{professormode}{
\begin{abstract}
The path to success in this class is to do many problems. Unlike other courses, exclusively doing reading(s) will not help. Coming to lecture is akin to watching workout videos; thinking about and solving problems on your own is the actual ``working out.''  Feel free to \qu{work out} with others; \textbf{I want you to work on this in groups.}

The problems below are color coded: \ingreen{green} problems are considered \textit{easy} and marked \qu{[easy]}; \inorange{yellow} problems are considered \textit{intermediate} and marked \qu{[harder]}, \inred{red} problems are considered \textit{difficult} and marked \qu{[difficult]} and \inpurple{purple} problems are extra credit. The \textit{easy} problems are intended to be ``giveaways'' if you went to class. Do as much as you can of the others; I expect you to at least attempt the \textit{difficult} problems. 

Bonus points are given as a bonus if the homework is typed using \LaTeX. Links to installing \LaTeX~and program for compiling \LaTeX~is found on the syllabus. You are encouraged to use \url{https://www.overleaf.com/}. If you are handing in homework this way, read the comments in the code; there is a line to comment out and you should replace my name with yours. The easiest way to use Overleaf is to make an Overleaf account, click the link to the \LaTeX~ template, go to \textit{File} and select \textit{Make a copy}. If you are asked to make drawings, you can take a picture of your handwritten drawing and insert them as figures or leave space using the \qu{$\backslash$vspace} command and draw them in after printing or attach them stapled.


\end{abstract}

\thispagestyle{empty}
\vspace{1cm}
NAME: \line(1,0){380}
\clearpage
}

\problem{These exercises will review the beginning of combinations.}


\begin{enumerate}

\easysubproblem{ What are the four assumptions we need in order to have a combination?}


\easysubproblem{Informally speaking, what is a combination?
}


\easysubproblem{Give three examples of a combination of 40 elements taken 5 at a time. Also provide four non-examples.}

\vspace{50mm}


\easysubproblem{Write down two symbols which denotes the number of combinations of $n$ distinct elements taken $k$ at a time.}

\newpage 

\end{enumerate}

\problem{These exercises will provide a more detailed perspective on combinations.}


\begin{enumerate}

\easysubproblem{Prove the following theorem: $\binom{n}{k}= \frac{P_{n,k}}{k!}$}

\vspace{150mm}

\easysubproblem{In the above question, $\binom{n}{k}$ is expressed in terms of permutations and factorials. Re-express $\binom{n}{k}$ by \textbf{only} using factorials. Hint: substitute in the formula for $P_{n,k}$. We also did this in class. }

\newpage

\easysubproblem{True or False? For a certain group of elements, if the order of selection does not matter, then we should select all those elements within one step.}

\vspace{10mm}

\intermediatesubproblem{Alice was asked the following question: From a group of three people, a president and two treasurers are to be chosen. How many outcomes are there? \\

Alice wrote down the following solution: 
\begin{tcolorbox}
We proceed in steps. 

\begin{enumerate}
    \item First choose a president. There are 3 choices.
    \item Then choose a subset of two treasurers. There are now $\binom{2}{2}$ choices.
\end{enumerate}

By The Generalized Principle of Counting, there are $3 \times \binom{2}{2} = 3$ choices.

\end{tcolorbox}

Draw a tree diagram for this experiment and use the diagram to explain why Alice's solution is correct. (That is, explain why we are allowed to invoke The Generalized Principle of Counting). }

\newpage 

\intermediatesubproblem{Alice was asked the following question: From a regular deck of 52 cards, suppose an outcome is a subset of 5 cards. How many outcomes are there where we have the denominations $a, a, a, b, b$ where $a$ and $b$ are distinct? For example, we obtain $\braces{A \clubsuit, A \diamondsuit, A \heartsuit, J \spadesuit, J \heartsuit} $.

Alice wrote down the following solution: 

\begin{tcolorbox}
We proceed in steps. 

\begin{enumerate}
    \item First choose a denomination for $a$. There are $\binom{13}{1}$ choices.
    \item Now choose three cards of that denomination. There are $\binom{4}{3}$ ways to do so.
    \item Then select a denomination for $b$. There are $\binom{12}{1}$ choices.
    \item Now choose two cards of this denomination. There are $\binom{4}{2}$ ways to do so.
\end{enumerate}

By The Generalized Principle of Counting, there are $ \binom{13}{1} \times \binom{4}{3} \times \binom{12}{1} \times \binom{4}{2}$ choices.

\end{tcolorbox}

Is Alice's solution correct? Explain why or why not. Hint: Look back at the question we considered in class regarding the full house.}

\newpage 

\hardsubproblem{Alice was asked the following question: From a regular deck of 52 cards, suppose an outcome is a subset of 6 cards. How many outcomes are there where we have the denominations $a, a, a, b, b, b$ where $a$ and $b$ are distinct? For example, we obtain $\braces{A \clubsuit, A \diamondsuit, A \heartsuit, J \spadesuit, J \heartsuit, J \diamondsuit} $.

Alice wrote down the following solution: 

\begin{tcolorbox}
We proceed in steps. 

\begin{enumerate}
    \item First choose a denomination for $a$. There are $\binom{13}{1}$ choices.
    \item Now choose three cards of that denomination. There are $\binom{4}{3}$ ways to do so.
    \item Then select a denomination for $b$. There are $\binom{12}{1}$ choices.
    \item Now choose three cards of this denomination. There are $\binom{4}{3}$ ways to do so.
\end{enumerate}

By The Generalized Principle of Counting, there are $ \binom{13}{1} \times \binom{4}{3} \times \binom{12}{1} \times \binom{4}{3}$ choices.

\end{tcolorbox}

Is Alice's solution correct? Explain why or why not.}

\newpage 

\easysubproblem{Prove that $\binom{n}{k} = \binom{n}{n-k}$.}


\vspace{60mm}

\easysubproblem{Intuitively explain why $\binom{25}{5} = \binom{25}{20}$.}

\vspace{60mm}

\easysubproblem{A committee of 7, consisting of 2 Republicans, 2 Democrats, and 3 Independents, is to be chosen from a group of 5 Republicans, 6 Democrats, and 4 Independents. How many committees are possible?} 

\newpage 

\hardsubproblem{ A dance class consists of 22 students, of which 10 are women and 12 are men. If 5 men and 5 women are to be chosen and then paired off, how many
results are possible? (Paired off means that each man is partnered to one woman as a dance partner).}

\vspace{80mm}



\intermediatesubproblem{Consider a group of 20 people. If everyone shakes hands with everyone else, then how many handshakes take place?} 

\vspace{60mm}

\easysubproblem{From a group of 5 women and 7 men, how many different committees are there which consist of 2 women and 3 men?}

\newpage 

\hardsubproblem{In the question above, suppose that one woman (Amber) is feuding with a man (Johnny) and the two of them refuse to serve together. How many different committees are possible? }

\end{enumerate}

\newpage 

\problem{We are almost finished with our study of combinatorial analysis. Our last step is to bridge the Chapter 2 material (counting) with the Chapter 1 material (probability). These set of exercises will begin to establish this connection. At the end of Chapter 1, we introduced the notion of simple sample spaces. We then stated a theorem that  said if $S$ is a simple sample space with $k$ elements, and $A$ is an event with $m$ elements, then 
\begin{align*}
    P(A) = \frac{m}{k} = \frac{\text{the number of elements in $A$}}{\text{the number of elements in $S$}}
\end{align*}
Take a few moments to remind yourself of why this result is true. It may be helpful for the next set of exercises! }


\begin{enumerate}

\easysubproblem{ An experiment involves flipping a fair coin four times, and recording the resulting sequence of heads and tails in the order they appear. Write down all of the outcomes in $S$. }

\vspace{60mm}

\easysubproblem{If the coin is truly fair, then the probability of each outcome in the above part should be 1/16. Let the events $A$, $B$, $C$, and $D$ be given by $A = \{\text{at least 3 heads} \}, B = \{\text{at most 2 heads} \}, C = \{ \text{heads on the third toss} \}$, and $D = \{\text{1 head and 3 tails} \}$.  Find $P(A)$. Hint: count the number of outcomes in $A$ and divide by the number of outcomes in $S$ which is 16.
}

\vspace{30mm}

\easysubproblem{Find $P(B)$.}

\vspace{30mm}

\easysubproblem{Find $P(A \cap B)$. Hint: count the number of outcomes in $A \cap B$ and divide by the number of outcomes in $S$ which is 16.  }

\vspace{30mm}

\easysubproblem{Find $P(A \cap C)$.}

\vspace{30mm}

\easysubproblem{Find $P(A \cup C)$.}

\vspace{30mm}

\easysubproblem{Find $P(B \cap D)$.}
\end{enumerate}
\newpage

\problem 
\begin{enumerate}
 
\easysubproblem{ \textbf{Optional}: This question will not be graded. If you can, try to explain the following meme: 
\\
\begin{center}
\includegraphics[width=7cm]{memes.png}
\end{center}}

\end{enumerate}

\end{document}
